\documentclass[a4paper, 12pt]{article}

% золотой набор преамблы
\usepackage[english, russian]{babel} % Подключает языковые модули. Настраивает правила переноса слов и автоматически переводит системные слова (например, вместо "Contents" пишет "Оглавление", вместо "Theorem" — "Теорема"). Последний язык в списке (russian) становится основным.

\usepackage[T2A]{fontenc} % Задает внутреннюю кодировку шрифтов LaTeX. T2A — это стандартная кодировка для поддержки кириллицы. Благодаря ей в итоговом PDF-файле русский текст будет нормально копироваться и искаться.

\usepackage[utf8]{inputenc} % Говорит компилятору, что твой исходный код .tex написан в кодировке UTF-8. Как я говорил ранее, в современном Overleaf этот пакет уже встроен по умолчанию, его можно опускать.

\usepackage{amsmath, amsfonts, amssymb, amsthm, mathtools} % интерфейс работы с теоремами обеспечивается пакетом amsthm
\usepackage{graphicx} % это главный пакет в LaTeX для вставки и редактирования изображений.
\usepackage[left=2cm, right=2cm, top=2cm, bottom=2cm]{geometry}
\usepackage{hyperref} % работа с ссылками, делает кликабельным
\hypersetup{
    colorlinks=true,    % Включает цветной текст ссылок вместо уродливых рамок
    linkcolor=blue,     % Цвет внутренних ссылок (оглавление, формулы)
    urlcolor=cyan,      % Цвет внешних веб-ссылок
    citecolor=green,    % Цвет ссылок на литературу
    pdftitle={Конспект по матанализу}, % Метаданные для PDF-файла
    pdfauthor={Студент 1-го курса}
}

\usepackage{indentfirst} % пакет для абзаца каждой секции
\usepackage{titleps} % пакент для колонтитулов
\usepackage{soulutf8} % пакет для разного вида шрифта
\usepackage{xcolor} % для работы с цветами

% для работы с таблицами
\usepackage{array}
\usepackage{tabularx}
\usepackage{multirow}

% стиль страницы
\newpagestyle{main}{ % main - название стиля, и в других фигурных скобках это содержимое
    \setheadrule{0.4pt} % толщина линии отделяющий хедер или футер от текста
    \sethead{ФПМИ Конспект}{}{Глава \thesection} % верхний колонтитул (header) обязательные три параметра!
    % \thesection — выведет номер раздела (например, 1.2)
    % \sectiontitle — автоматически выведет имя текущего раздела
    \setfootrule{0.4pt}  % толщина линии отделяющий хедер или футер от текста
    \setfoot{Конспект лекций}{\thepage}{2026} % нижний колонтитул (footer)
}
\pagestyle{main} % команда для приминения стиля страницы, можно использовать в каком-то определенном промежутке кода

% стиль plain
\theoremstyle{plain} % модификатор, который задает стиль всем следуюшим теоремам, пока не переопределить.
\newtheorem{theorem}{Теорема}
\newtheorem{lemma}{Лемма}
\newtheorem{axeom}{Аксиома}
\newtheorem{corollary}{Следствие}[theorem] % последний необязательный аргумент, говорит о нумерации следствия данной теоремы.

% стиль defenition
\theoremstyle{definition}
\newtheorem{definition}{Определение}[theorem] % привязанность к данной теореме
\newtheorem{property}{Свойство}[theorem]
\newtheorem{exercise}{Задача}[section]

% стиль remark
\theoremstyle{remark}
\newtheorem{remark}[theorem]{Замечание}